\chapter{Krav og akseptanse}
\label{ch:requirements}
\section{Funksjonelle krav}
De opprinnelige krav-ID-ene beholdes. Formuleringene nedenfor presiserer
hvordan kravene kan testes, uten å sette likhetstegn mellom kode og fysisk
oppnådd funksjon. \emph{Implementert} betyr at funksjonen finnes i den
aktive koden; \emph{godkjent} krever også dokumenterte fysiske resultater.

\begin{tabularx}{\textwidth}{@{}P{12mm}Y Y@{}}
\toprule ID & Etterprøvbart krav & Status og akseptanse\\\midrule
FR1 & Kjøre autonomt etter bevisst start. & A-gate og løpende styring testet på vert. Autonom løype gjenstår.\\
FR2 & Oppdage hindring og ukjent sensormåling. & Avstandsgrenser og ugyldige data testet. Tre fysiske kanaler må kalibreres.\\
FR3 & Velge retning fra sensordata. & Første gjennomførbare skannekandidat testet; fysiske skannevinkler gjenstår.\\
FR4 & Utføre unnamanøver ved hindring. & Begge svinggrener og avbrudd testet mot fakes. Chassisklaring må måles.\\
FR5 & Styre fire Mecanum-motorer korrekt. & Motorkommandoer bygger. Alle hjulfortegn og begge svingretninger må bekreftes løftet.\\
FR6 & Fortsette etter unnamanøver. & Sving går via ny verifikasjon. Sammenhengende fysisk sekvens gjenstår.\\\bottomrule
\end{tabularx}

\section{Sanntids- og systemkrav}
RT1 krever regelmessig sensorovervåking, RT2 reaksjon innen en begrunnet
frist og RT3 at kritiske aktiviteter møter sine deadlines. Disse kravene
er ennå ikke godkjent. Koden konfigurerer minst 65 ms mellom akustiske
triggere, men en periodeparameter er ikke en målt maksimal reaksjonstid.
Kapittel~\ref{ch:timing} beskriver analyse og manglende måleverdier.

SR1--SR4 gjelder micro:bit V2, Ada, Mecanum-bil med fire motorer og
tilkoblede sensorer. Brukeren har bekreftet oppsettet; byggene retter seg
mot \kode{embedded-nrf52833}. Motor- og sensorforsyning er frakoblet i
siste dokumenterte fysiske oppsett. Forsyning, ECHO-nivåer og faktisk
motorrespons må verifiseres før integrasjonstesten kan godkjennes.

\section{Utledede feil- og brukerkrav}
\begin{tabularx}{\textwidth}{@{}P{13mm}Y Y@{}}\toprule
ID & Krav & Kontrollerbar egenskap\\\midrule
S1 & Ingen automatisk kjøring ved oppstart. & A må forespørre start; initialisering alene gir ingen kjørekommando.\\
S2 & B har forrang over A. & Stopp er latchet; frigivelse av B eller et nytt A-trykk opphever ikke stoppen.\\
S3 & Feil må ikke føre til automatisk omstart. & Første feil lagres; én stoppoperasjon; deretter parkert feiltilstand.\\
S4 & Ikke hevde fysisk stopp fra busskvittering. & Programmet skiller sendt stopp fra ubekreftet stopp.\\
S5 & Gamle skannedata skal ikke tillate sving. & Eldste stråles alder kontrolleres også etter siste sensorkontroll før kjøring.\\
S6 & Lys må ikke konkurrere om motorbussen. & Én GPIO-eier; protected object deler kun status.\\\bottomrule
\end{tabularx}

\section{Faglig sporbarhet}
Sterk typing og pakker knyttes til \kode{Adaptive_Policy}; exceptions og
propagation til \kode{Adaptive_Errors} og \kode{Adaptive_Fault_Handler};
tasking, delte ressurser og prioritet til \kode{Adaptive_Lights} og
\kode{Adaptive_Lamp_Control}. Responstid og blokkering behandles mot
den faktiske sekvensielle styringen. Kontrakter og formell analyse knyttes
til den reelle lyspolicyen, og testbarhet til den generiske kontrolleren.

Akseptansetesten skal registrere firmwarehash, oppsett, rådata, kriterium,
utfall og alle mislykkede forsøk. Et krav lukkes først når kriteriet er
bestått med den samme konfigurasjonen som demonstreres. Kravene om
fysisk funksjon og tid kan derfor ikke lukkes av vertsbaserte tester.
